% GENERATED FILE. Edit canonical lessons, then run npm run generate:books.
% canonical-source-hash: fnv1a64:7e0213f42b5b6627
% canonical-lessons: SA-C525-mallah, SA-C525-bahuyuddham, SA-C525-patayati, SA-C525-uttolayati, SA-C525-ahvanam, SA-R525-review-518-525

\chapter{In the Wrestling Pit}
\label{ch:sa-in-the-wrestling-pit}
\hlchaptermodality{525}

\begin{chapteropening}
\textbf{By the end of this chapter:} \emph{I can describe a wrestling match.}

Complete the last lesson of chapter 525: I can describe a wrestling match.
\end{chapteropening}

\section[\texorpdfstring{\sk{मल्लः} (mallaḥ)}{मल्लः (mallaḥ)}]{\sk{मल्लः} (mallaḥ) — a wrestler}

\label{lesson:SA-C525-mallah}



\begin{quote}
Before the new one: say the Sanskrit for beaten, then the Sanskrit for hits.
\end{quote}

\subsection*{You'll want to know: \sk{मल्लः}}
\textbf{\sk{मल्लः}} — \emph{mallaḥ} — \textquotedblleft{}a wrestler\textquotedblright{}.

\begin{quote}
\sk{ग्रामस्य} \sk{मल्लः} \sk{बलवान्} \sk{अस्ति।}
\end{quote}

\begin{quote}
\emph{grāmasya mallaḥ balavān asti.}
\end{quote}

\begin{quote}
The village wrestler is strong.
\end{quote}

A masculine noun.

\begin{cousinweb}
\textbf{\sk{मल्लः}} (\emph{mallaḥ}) is the word MW records for a wrestler, a boxer; MW gives no derivation, so its origin is uncertain (MW).
\end{cousinweb}

\begin{grammarlens}[title={What you've built}]
One more word for describing what you lived and felt.
\end{grammarlens}

\subsection*{Your turn}
\begin{itemize}
\raggedright
  \item \emph{Say it:} \emph{mallaḥ}
  \item \emph{Say it:} \emph{mallaḥ}, once more
  \item \emph{Say it:} \emph{grāmasya mallaḥ balavān asti.} — The village wrestler is strong.
  \item \emph{Recall:} say the Sanskrit for beaten, then the Sanskrit for hits, then say \emph{mallaḥ} again
\end{itemize}

\begin{tcolorbox}[breakable,colback=teal!4,colframe=teal!35!black,title={Before you move on}]
What does \sk{मल्लः} mean? (\textquotedblleft{}A wrestler\textquotedblright{}.) Say it once more.
\end{tcolorbox}
\section[\texorpdfstring{\sk{बाहुयुद्धम्} (bāhuyuddham)}{बाहुयुद्धम् (bāhuyuddham)}]{\sk{बाहुयुद्धम्} (bāhuyuddham) — wrestling}

\label{lesson:SA-C525-bahuyuddham}



\begin{quote}
Before the new one: say the Sanskrit for hits, then the Sanskrit for a wrestler.
\end{quote}

\subsection*{You'll want to know: \sk{बाहुयुद्धम्}}
\textbf{\sk{बाहुयुद्धम्}} — \emph{bāhuyuddham} — \textquotedblleft{}wrestling\textquotedblright{}.

\begin{quote}
\sk{सायं} \sk{ग्रामे} \sk{बाहुयुद्धं} \sk{भवति।}
\end{quote}

\begin{quote}
\emph{sāyaṃ grāme bāhuyuddhaṃ bhavati.}
\end{quote}

\begin{quote}
There is wrestling in the village in the evening.
\end{quote}

A neuter noun.

\begin{cousinweb}
\textbf{\sk{बाहुयुद्धम्}} (\emph{bāhuyuddham}) is built from \emph{bāhu}, \textquotedblleft{}an arm\textquotedblright{}, and \emph{yuddha}, \textquotedblleft{}a fight\textquotedblright{}; MW gives \emph{bāhuyuddha} as an arm-fight, a close fight (MW).
\end{cousinweb}

\begin{grammarlens}[title={What you've built}]
One more word for describing what you lived and felt.
\end{grammarlens}

\subsection*{Your turn}
\begin{itemize}
\raggedright
  \item \emph{Say it:} \emph{bāhuyuddham}
  \item \emph{Say it:} \emph{bāhuyuddham}, once more
  \item \emph{Say it:} \emph{sāyaṃ grāme bāhuyuddhaṃ bhavati.} — There is wrestling in the village in the evening.
  \item \emph{Recall:} say the Sanskrit for hits, then the Sanskrit for a wrestler, then say \emph{bāhuyuddham} again
\end{itemize}

\begin{tcolorbox}[breakable,colback=teal!4,colframe=teal!35!black,title={Before you move on}]
What does \sk{बाहुयुद्धम्} mean? (\textquotedblleft{}Wrestling\textquotedblright{}.) Say it once more.
\end{tcolorbox}
\section[\texorpdfstring{\sk{पातयति} (pātayati)}{पातयति (pātayati)}]{\sk{पातयति} (pātayati) — throws down}

\label{lesson:SA-C525-patayati}



\begin{quote}
Before the new one: say the Sanskrit for a wrestler, then the Sanskrit for wrestling.
\end{quote}

\subsection*{You'll want to know: \sk{पातयति}}
\textbf{\sk{पातयति}} — \emph{pātayati} — \textquotedblleft{}throws down\textquotedblright{}.

\begin{quote}
\sk{एकः} \sk{मल्लः} \sk{अन्यं} \sk{भूमौ} \sk{पातयति।}
\end{quote}

\begin{quote}
\emph{ekaḥ mallaḥ anyaṃ bhūmau pātayati.}
\end{quote}

\begin{quote}
One wrestler throws the other to the ground.
\end{quote}

Past (he …): \sk{सः} \sk{मल्लं} \sk{भूमौ} \sk{पातितवान्}, \emph{saḥ mallaṃ bhūmau pātitavān}.

\begin{cousinweb}
\textbf{\sk{पातयति}} (\emph{pātayati}) comes from the causative of the root √pat, \textquotedblleft{}to fall\textquotedblright{}, which also gives \emph{patati}: to make fall, to throw down (MW).
\end{cousinweb}

\begin{grammarlens}[title={What you've built}]
One more word for describing what you lived and felt.
\end{grammarlens}

\subsection*{Your turn}
\begin{itemize}
\raggedright
  \item \emph{Say it:} \emph{pātayati}
  \item \emph{Say it:} \emph{pātayati}, once more
  \item \emph{Say it:} \emph{ekaḥ mallaḥ anyaṃ bhūmau pātayati.} — One wrestler throws the other to the ground.
  \item \emph{Recall:} say the Sanskrit for a wrestler, then the Sanskrit for wrestling, then say \emph{pātayati} again
\end{itemize}

\begin{tcolorbox}[breakable,colback=teal!4,colframe=teal!35!black,title={Before you move on}]
What does \sk{पातयति} mean? (\textquotedblleft{}Throws down\textquotedblright{}.) Say it once more.
\end{tcolorbox}
\section[\texorpdfstring{\sk{उत्तोलयति} (uttolayati)}{उत्तोलयति (uttolayati)}]{\sk{उत्तोलयति} (uttolayati) — hoists; lifts up}

\label{lesson:SA-C525-uttolayati}



\begin{quote}
Before the new one: say the Sanskrit for wrestling, then the Sanskrit for throws down.
\end{quote}

\subsection*{You'll want to know: \sk{उत्तोलयति}}
\textbf{\sk{उत्तोलयति}} — \emph{uttolayati} — \textquotedblleft{}hoists; lifts up\textquotedblright{}.

\begin{quote}
\sk{मल्लः} \sk{गुरुं} \sk{भारम्} \sk{उत्तोलयति।}
\end{quote}

\begin{quote}
\emph{mallaḥ guruṃ bhāram uttolayati.}
\end{quote}

\begin{quote}
The wrestler lifts a heavy load.
\end{quote}

Past (he …): \sk{सः} \sk{भारम्} \sk{उत्तोलितवान्}, \emph{saḥ bhāram uttolitavān}.

\begin{cousinweb}
\textbf{\sk{उत्तोलयति}} (\emph{uttolayati}) is built on \emph{ud}, \textquotedblleft{}up\textquotedblright{}, and the causative of the root √tul, \textquotedblleft{}to lift, to weigh\textquotedblright{}, which also gives \emph{tolayati}; MW gives \emph{uttolayati} as to take up, to raise up (MW).
\end{cousinweb}

\begin{grammarlens}[title={What you've built}]
One more word for describing what you lived and felt.
\end{grammarlens}

\subsection*{Your turn}
\begin{itemize}
\raggedright
  \item \emph{Say it:} \emph{uttolayati}
  \item \emph{Say it:} \emph{uttolayati}, once more
  \item \emph{Say it:} \emph{mallaḥ guruṃ bhāram uttolayati.} — The wrestler lifts a heavy load.
  \item \emph{Recall:} say the Sanskrit for wrestling, then the Sanskrit for throws down, then say \emph{uttolayati} again
\end{itemize}

\begin{tcolorbox}[breakable,colback=teal!4,colframe=teal!35!black,title={Before you move on}]
What does \sk{उत्तोलयति} mean? (\textquotedblleft{}Hoists; lifts up\textquotedblright{}.) Say it once more.
\end{tcolorbox}
\section[\texorpdfstring{\sk{आह्वानम्} (āhvānam)}{आह्वानम् (āhvānam)}]{\sk{आह्वानम्} (āhvānam) — a challenge; a summons}

\label{lesson:SA-C525-ahvanam}



\begin{quote}
Before the new one: say the Sanskrit for throws down, then the Sanskrit for hoists.
\end{quote}

\subsection*{You'll want to know: \sk{आह्वानम्}}
\textbf{\sk{आह्वानम्}} — \emph{āhvānam} — \textquotedblleft{}a challenge; a summons\textquotedblright{}.

\begin{quote}
\sk{मल्लः} \sk{आह्वानं} \sk{स्वीकरोति।}
\end{quote}

\begin{quote}
\emph{mallaḥ āhvānaṃ svīkaroti.}
\end{quote}

\begin{quote}
The wrestler accepts the challenge.
\end{quote}

A neuter noun.

\begin{cousinweb}
\textbf{\sk{आह्वानम्}} (\emph{āhvānam}) is built on \emph{ā} and the root √hve, \textquotedblleft{}to call\textquotedblright{}, which also gives \emph{āhvayati}; MW gives \emph{āhvāna} as calling, a summons, and a challenge (MW).
\end{cousinweb}

\begin{grammarlens}[title={What you've built}]
One more word for describing what you lived and felt.
\end{grammarlens}

\subsection*{Your turn}
\begin{itemize}
\raggedright
  \item \emph{Say it:} \emph{āhvānam}
  \item \emph{Say it:} \emph{āhvānam}, once more
  \item \emph{Say it:} \emph{mallaḥ āhvānaṃ svīkaroti.} — The wrestler accepts the challenge.
  \item \emph{Recall:} say the Sanskrit for throws down, then the Sanskrit for hoists, then say \emph{āhvānam} again
\end{itemize}

\begin{tcolorbox}[breakable,colback=teal!4,colframe=teal!35!black,title={Before you move on}]
What does \sk{आह्वानम्} mean? (\textquotedblleft{}A challenge; a summons\textquotedblright{}.) Say it once more.
\end{tcolorbox}
\section[(dialogue)]{Looking back: from Colours on the Wall to In the Wrestling Pit}

\label{lesson:SA-R525-review-518-525}



\begin{quote}
Say the Sanskrit for a paintbrush and paint.
\end{quote}

\subsection*{Your turn}
Say these aloud:

\begin{itemize}
\raggedright
  \item a paintbrush, paint, a line, paints, a portrait
  \item a statue, carves, a craft, a creation, fashions
  \item a musical note, a raga, a tempo, singing, plays an instrument
  \item a musician, a string of an instrument, a hand drum, a tambourine, a resonant sound
  \item a dance step, spins round, a small bell, imitates, clapping
  \item an actor, acts a part, a jester, backstage, acting
  \item a contest, hurls, a victor, beaten, hits
  \item a wrestler, wrestling, throws down, hoists, a challenge
\end{itemize}

\begin{tcolorbox}[breakable,colback=teal!4,colframe=teal!35!black,title={Before you move on}]
A paintbrush? (\textbf{\sk{तूलिका}}, \emph{tūlikā}.) Paint? (\textbf{\sk{रंगः}}, \emph{raṅgaḥ}.) A line? (\textbf{\sk{रेखा}}, \emph{rekhā}.) Paints? (\textbf{\sk{आलिखति}}, \emph{ālikhati}.) A portrait? (\textbf{\sk{प्रतिकृतिः}}, \emph{pratikṛtiḥ}.) A statue? (\textbf{\sk{मूर्तिः}}, \emph{mūrtiḥ}.) Carves? (\textbf{\sk{उत्किरति}}, \emph{utkirati}.) A craft? (\textbf{\sk{शिल्पम्}}, \emph{śilpam}.) A creation? (\textbf{\sk{कृतिः}}, \emph{kṛtiḥ}.) Fashions? (\textbf{\sk{निर्माति}}, \emph{nirmāti}.) A musical note? (\textbf{\sk{स्वरः}}, \emph{svaraḥ}.) A raga? (\textbf{\sk{रागः}}, \emph{rāgaḥ}.) A tempo? (\textbf{\sk{लयः}}, \emph{layaḥ}.) Singing? (\textbf{\sk{गानम्}}, \emph{gānam}.) Plays an instrument? (\textbf{\sk{वादयति}}, \emph{vādayati}.) A musician? (\textbf{\sk{वादकः}}, \emph{vādakaḥ}.) A string of an instrument? (\textbf{\sk{तन्त्री}}, \emph{tantrī}.) A hand drum? (\textbf{\sk{डमरुः}}, \emph{ḍamaruḥ}.) A tambourine? (\textbf{\sk{मुरजः}}, \emph{murajaḥ}.) A resonant sound? (\textbf{\sk{नादः}}, \emph{nādaḥ}.) A dance step? (\textbf{\sk{पदन्यासः}}, \emph{padanyāsaḥ}.) Spins round? (\textbf{\sk{आवर्तते}}, \emph{āvartate}.) A small bell? (\textbf{\sk{किंकिणी}}, \emph{kiṅkiṇī}.) Imitates? (\textbf{\sk{अनुकरोति}}, \emph{anukaroti}.) Clapping? (\textbf{\sk{करतालिका}}, \emph{karatālikā}.) An actor? (\textbf{\sk{नटः}}, \emph{naṭaḥ}.) Acts a part? (\textbf{\sk{अभिनयति}}, \emph{abhinayati}.) A jester? (\textbf{\sk{विदूषकः}}, \emph{vidūṣakaḥ}.) Backstage? (\textbf{\sk{नेपथ्यम्}}, \emph{nepathyam}.) Acting? (\textbf{\sk{अभिनयः}}, \emph{abhinayaḥ}.) A contest? (\textbf{\sk{स्पर्धा}}, \emph{spardhā}.) Hurls? (\textbf{\sk{प्रक्षिपति}}, \emph{prakṣipati}.) A victor? (\textbf{\sk{विजेता}}, \emph{vijetā}.) Beaten? (\textbf{\sk{पराजितः}}, \emph{parājitaḥ}.) Hits? (\textbf{\sk{प्रहरति}}, \emph{praharati}.) A wrestler? (\textbf{\sk{मल्लः}}, \emph{mallaḥ}.) Wrestling? (\textbf{\sk{बाहुयुद्धम्}}, \emph{bāhuyuddham}.) Throws down? (\textbf{\sk{पातयति}}, \emph{pātayati}.) Hoists? (\textbf{\sk{उत्तोलयति}}, \emph{uttolayati}.) A challenge? (\textbf{\sk{आह्वानम्}}, \emph{āhvānam}.)
\end{tcolorbox}
